\chapter{Solution Design}
\label{cap:design}

This chapter presents the design of TalentHub: the high-level container
view, the onion architecture that organises the code, the domain model and
its relational schema, the job-anchored matching algorithm, the AI pipelines
that ingest unstructured text, and the security, compliance and frontend
decisions. It closes with a decision log that records the main trade-offs.

%----------------------------------------------------------------------------------------

\section{Architectural Overview}
\label{sec:design-overview}

At the highest level the system is a single web application supported by
managed services and one specialised sidecar. The \textbf{Next.js
application} hosts both the user interface and the server-side API routes,
and is the only component a user's browser talks to directly. It depends on
three \textbf{Supabase} managed services --- PostgreSQL for data,
Authentication for Google OAuth, and Storage for uploaded CVs --- and on two
external \textbf{Large Language Model} providers --- a primary and a
fallback --- for parsing. The conversational \textbf{AI Interviewer} is a
separate Python \textit{FastAPI} service: the browser opens an interview
popup that communicates with the Next.js API routes, which in turn delegate
the conversation to the sidecar through a signed-token contract.
\tabref{tab:containers} summarises the containers and their responsibilities.

\begin{table}[ht]
\centering
\caption[System containers]{Top-level containers and their responsibilities.}
\label{tab:containers}
\small
\begin{tabular}{p{0.26\linewidth}p{0.30\linewidth}p{0.32\linewidth}}
\toprule
\textbf{Container} & \textbf{Technology} & \textbf{Responsibility} \\
\midrule
Web application & Next.js 16 / TypeScript on Vercel & UI, API routes, use cases, repositories \\
Database & Supabase PostgreSQL & Persistent relational store \\
Authentication & Supabase Auth (Google OAuth) & Identity, session, role claim \\
Object storage & Supabase Storage & Uploaded CV files \\
LLM providers & Groq (primary), OpenAI (fallback) & CV and job-description parsing \\
AI Interviewer & Python FastAPI sidecar & Conversational assessment + CEFR evaluation \\
Email & Resend & HR-to-candidate outreach; live delivery gated on sender-domain verification \\
\bottomrule
\end{tabular}
\end{table}

\figref{fig:c4-container} renders the same containers as a context view,
making explicit which actor talks to which component and where each external
dependency sits. The single inward-facing fact it captures is that the
browser only ever reaches the Next.js application; every managed service and
external provider is contacted server-side, which is the structural
precondition for the application-edge authorisation model defended in
\secref{sec:design-security}.

\begin{figure}[ht]
\centering
\begin{tikzpicture}[
  font=\scriptsize,
  person/.style={draw, rounded corners, fill=lstkeyword!12, text width=2.0cm, align=center, minimum height=0.9cm},
  core/.style={draw, thick, fill=lstkeyword!30, text width=3.0cm, align=center, minimum height=1.1cm},
  ext/.style={draw, fill=lstcomment!18, text width=2.2cm, align=center, minimum height=0.9cm}]
% actors
\node[person] (hr) at (0,0) {HR recruiter (browser)};
\node[person] (cand) at (4.2,0) {Candidate (browser)};
% core
\node[core] (web) at (2.1,-2.0) {Next.js web app: UI + API routes + use cases (Vercel)};
% external left/right
\node[ext] (llm) at (-2.6,-2.0) {LLM providers Groq / OpenAI};
\node[ext] (ai) at (6.8,-2.0) {FastAPI AI Interviewer};
% managed services
\node[ext] (db) at (-1.2,-4.2) {Supabase PostgreSQL};
\node[ext] (auth) at (2.1,-4.2) {Supabase Auth (OAuth)};
\node[ext] (store) at (5.4,-4.2) {Supabase Storage};
% relations
\foreach \a in {hr,cand} \draw[-{Stealth}] (\a) -- (web);
\foreach \b in {llm,ai,db,auth,store} \draw[-{Stealth}, lstcomment] (web) -- (\b);
\end{tikzpicture}
\caption[Container view]{Container (context) view of TalentHub.}
\label{fig:c4-container}
\end{figure}

%----------------------------------------------------------------------------------------

\section{Onion / Clean Architecture}
\label{sec:design-onion}

The code follows an \gls{onion}, chosen and enforced from the start
\parencite{PalermoOnion2008,Martin2017}. The
guiding rule is that \emph{dependencies always point inward}: the Domain at
the centre has zero external dependencies, the Application layer orchestrates
the Domain, the Infrastructure layer implements the Domain's port interfaces,
and the Presentation layer is a thin shell over the Application. Crucially,
Infrastructure \emph{depends on} the Domain (by implementing its ports) while
the Domain never imports Infrastructure:

\begin{quote}
Presentation $\rightarrow$ Application $\rightarrow$ Domain $\leftarrow$
Infrastructure
\end{quote}

\figref{fig:onion} draws this as concentric layers: the dependency-free
Domain at the centre, the Application layer wrapping it, and Infrastructure
and Presentation forming the outer shell.

\begin{figure}[ht]
\centering
\begin{tikzpicture}[font=\footnotesize]
\filldraw[fill=lstkeyword!8, draw=lstkeyword, thick, rounded corners] (-5.0,-3.4) rectangle (5.0,3.4);
\node[anchor=north] at (0,3.25) {\textbf{Presentation} --- API routes, React components};
\filldraw[fill=lstkeyword!18, draw=lstkeyword, thick, rounded corners] (-4.0,-2.75) rectangle (4.0,2.75);
\node[anchor=north] at (0,2.6) {\textbf{Infrastructure} --- Supabase repos, LLM, storage, email};
\filldraw[fill=lstkeyword!32, draw=lstkeyword, thick, rounded corners] (-2.9,-2.0) rectangle (2.9,2.0);
\node[anchor=north] at (0,1.85) {\textbf{Application} --- use cases, DTOs};
\filldraw[fill=lstkeyword!55, draw=lstkeyword, thick, rounded corners] (-1.7,-1.0) rectangle (1.7,1.0);
\node[align=center] at (0,0) {\textbf{Domain}\\\scriptsize ports, value objects,\\\scriptsize pure services};
\draw[-{Stealth}, very thick, lststring] (4.6,-3.0) -- (1.9,-0.7);
\node[lststring, font=\scriptsize, anchor=west] at (3.0,-3.15) {dependencies point inward};
\end{tikzpicture}
\caption[Onion architecture]{The onion architecture: a dependency-free Domain
at the core, with Application, Infrastructure and Presentation around it.}
\label{fig:onion}
\end{figure}

\begin{description}
	\item[Domain (\texttt{src/server/domain/}).] Pure business logic ---
	value objects (CEFR levels, scoring weights, status configuration,
	thresholds), pure domain services (scoring, matching), and the port
	interfaces that abstract persistence and external services. No
	database, HTTP or framework code.
	\item[Application (\texttt{src/server/application/}).] Use cases that
	orchestrate domain logic and ports, plus DTOs validated with Zod. It
	never imports a concrete repository or external SDK.
	\item[Infrastructure (\texttt{src/server/infrastructure/}).] The only
	layer that knows about external services: Supabase repositories, the
	LLM parsing services, storage, email and logging --- each implementing
	a domain port.
	\item[Presentation (\texttt{src/app/}, \texttt{src/client/}).] Next.js
	API routes and React components. Thin: it parses the request, calls a
	use case, and renders the result.
\end{description}

This discipline was chosen for testability (Domain and Application are
tested without mocks of infrastructure), swappability (the migration from an
\gls{orm}-and-Neon stack to Supabase required only new repository implementations,
with no domain or use-case changes), and provider flexibility (changing LLM
provider is a new service implementation behind an existing port). The
alternatives --- MVC, a permissive three-tier split, microservices, or a
bag of serverless functions --- were rejected as either too coupled or too
heavy for a two-person team. This same separation is what later allowed the
AI Interviewer to be extracted into its own process behind an unchanged
contract.

\subsection{Dependency Injection and the Composition Root}
\label{sec:design-di}

The inward dependency rule is realised in practice through constructor
injection and a single composition root. Each use-case class receives its
collaborators --- always as port \emph{interfaces}, never as concrete classes
--- through its constructor, so the use case is unaware of which technology
fulfils the contract. A use case is therefore constructed against
\texttt{ICandidateRepository}, not against \texttt{SupabaseCandidateRepository},
and the same use case could be exercised against an in-memory fake in a unit
test without changing a line of its code.

The wiring itself is centralised. A single container module instantiates the
concrete infrastructure implementations and a factory module imports those
instances to assemble pre-wired use cases; these two files are the only place
in the codebase that names a concrete repository or service. The arrangement
also expresses the Interface Segregation Principle directly: each use case
receives only the collaborators it genuinely needs ---
\texttt{CandidateUseCases} a single candidate repository,
\texttt{UploadUseCases} the six its pipeline requires (candidate repository,
CV parser, storage, text extraction, deduplication and parsing-job
repositories) --- and nothing more, which contains the impact of any
infrastructure change and keeps the dependency graph legible from the
constructor signatures alone.

%----------------------------------------------------------------------------------------

\section{Domain Model}
\label{sec:design-domain}

The domain is inherently relational and is centred on two aggregates
\parencite{Evans2003,Vernon2013IDDD}. A
\textbf{Candidate} owns the facts derived from a CV --- \textit{Experiences}
(each tagged with one or more fields of work), \textit{Education} entries,
\textit{Languages} with CEFR levels, \textit{Skills}, \textit{Tags} and
\textit{Notes} --- and accumulates process artefacts: \textit{Assessments}
and their \textit{Results}, \textit{Improvement Tracks}, \textit{Job
Matches}, \textit{Job Applications} and \textit{Notifications}. A
\textbf{Job} owns its \textit{Applications}, cached \textit{Matches},
\textit{Assessment Templates} and \textit{Shortlist} entries. Around these
sit supporting entities --- \textit{Promo Campaigns} (which generate
notifications), \textit{Parsing Jobs} (which track bulk uploads),
\textit{Ambassador Programs} with their applications, and \textit{Candidate
Segments} --- giving thirty-three application tables in total.
The business rules that drive scoring and assessment are not scattered
through the codebase but are encoded as immutable value objects
\parencite{Fowler2002PoEAA} in the Domain
layer, so that changing a weight or a threshold is a single, reviewable edit
in one file. \tabref{tab:value-objects} lists the principal constants and
the rules they encode. Expressing these as domain constants rather than as
hard-coded numeric values scattered through the codebase is what makes the scoring ``transparent'' in
the sense promised to the client: every component of a score traces back to a
named constant whose value is visible and version-controlled.

\begin{table}[ht]
\centering
\caption[Principal domain value objects]{Principal domain value objects in
the Domain layer.}
\label{tab:value-objects}
\small
\begin{tabular}{p{0.30\linewidth}p{0.60\linewidth}}
\toprule
\textbf{Constant} & \textbf{Encoded rule} \\
\midrule
\texttt{CEFR\_LEVELS} & The six-level proficiency scale A1--C2 \\
\texttt{CV\_SCORING\_WEIGHTS} & Language 35\%, experience relevance 25\%, education 15\%, location 15\%, years of experience 10\% \\
\texttt{EDUCATION\_LEVEL\_SCORES} & High school 20, vocational 40, bachelor 60, master 80, PhD 100 \\
\texttt{ASSESSMENT\_DEFAULT\_WEIGHTS} & General-assessment rubric: grammar, vocabulary, clarity, fluency and customer handling at 20\% each \\
\texttt{BORDERLINE\_THRESHOLD} & Borderline band: score between 45 and 60 \\
\texttt{MAGIC\_LINK\_EXPIRY\_HOURS} & Assessment-link validity of 48 hours \\
\texttt{MAX\_FILE\_SIZE\_MB} / \texttt{MAX\_BULK\_FILES} & Upload limits: 10\,MB per file, 500 files per batch \\
\bottomrule
\end{tabular}
\end{table}

\texttt{ASSESSMENT\_DEFAULT\_WEIGHTS} applies to the general assessment
template --- the customer-service and technical rubrics --- and is distinct
from the language interview, which instead averages the three \gls{cefr}
sub-scores at equal weight (\secref{sec:impl-scoring}).

The candidate itself carries a lifecycle status that the domain treats as a
state machine rather than a free-text label. \figref{fig:status-sm} shows
the permitted transitions: a candidate enters as \texttt{NEW}, becomes
\texttt{PARSED} once the CV pipeline completes, is \texttt{SCREENED} after
matching or assessment, and then moves to an outcome. A borderline result is
not a dead end --- it routes the candidate onto an improvement track and back
for re-evaluation, which is the mechanism behind the develop-and-retain pillar
of the framework in \chapref{cap:analysis}. Encoding these as named states
with explicit transitions keeps the status filters in the talent pool, the
notifications fired on change, and the analytics buckets all consistent with a
single source of truth.

\begin{figure}[ht]
\centering
\resizebox{\textwidth}{!}{%
\begin{tikzpicture}[font=\scriptsize, >=Stealth, node distance=6mm,
  st/.style={draw, rounded corners, fill=lstkeyword!18, align=center, minimum height=0.85cm, text width=2.0cm},
  to/.style={-{Stealth}, draw, line width=0.5pt},
  lbl/.style={font=\scriptsize, inner sep=1.5pt, fill=white, fill opacity=0.85, text opacity=1}]
% --- states -----------------------------------------------------------------
\node[st] (new) at (0,0)     {NEW};
\node[st] (par) at (3.0,0)   {PARSED};
\node[st] (scr) at (6.0,0)   {SCREENED};
\node[st] (off) at (9.4,0)   {OFFER\_SENT};
\node[st] (hir) at (12.6,0)  {HIRED};
\node[st] (bor) at (6.0,-2.7)  {BORDERLINE};
\node[st] (imp) at (9.4,-2.7)  {ON\_IMPROV.\\TRACK};
\node[st] (rej) at (7.7,-5.4)  {REJECTED};
% --- initial pseudo-state ---------------------------------------------------
\fill (-1.5,0) circle (2.5pt);
\draw[to] (-1.5,0) -- (new);
% --- happy path (top row) ---------------------------------------------------
\draw[to] (new) -- (par);
\draw[to] (par) -- (scr);
\draw[to] (scr) -- node[lbl,above]{pass} (off);
\draw[to] (off) -- node[lbl,above]{accepts} (hir);
% --- borderline / improvement loop ------------------------------------------
\draw[to] (scr) -- node[lbl,right]{borderline} (bor);
\draw[to] (bor) -- node[lbl,above]{assign} (imp);
\draw[to] (imp.north) to[out=120, in=-25]
      node[lbl,pos=0.55,above]{re-evaluate} (scr.south east);
% --- rejection transitions (outer lanes, clear of the middle boxes) ---------
\draw[to] (scr.south west) to[out=-140, in=165, looseness=1.15]
      node[lbl,pos=0.5,left]{fail} (rej.west);
\draw[to] (off.south east) to[out=-40, in=15, looseness=1.15]
      node[lbl,pos=0.5,right]{declines} (rej.east);
\end{tikzpicture}}
\caption[Candidate lifecycle state machine]{Candidate lifecycle as a state
machine.}
\label{fig:status-sm}
\end{figure}
%----------------------------------------------------------------------------------------

\section{Database Design}
\label{sec:design-database}

The store is a Supabase-managed PostgreSQL database accessed directly
through the Supabase client, with \textbf{no ORM}. The schema lives in a
single canonical SQL file (\texttt{00000000000000\_schema.sql}) into which
every per-feature change was consolidated, so a fresh database is created by
running one script.

\paragraph{Conventions.}
Database columns are \texttt{snake\_case}; JavaScript objects are
\texttt{camelCase}; conversion happens at the repository boundary through
\texttt{camelizeKeys} and \texttt{snakeifyKeys}. Fields stored as JSONB ---
parsed CV data, evaluation rationales, telemetry --- are listed in a
\texttt{JSONB\_KEYS} opt-out so that their internal keys are \emph{not}
recursively camelised and survive round-trips verbatim. Primary keys are
\texttt{text} values generated at the application layer with
\texttt{crypto.randomUUID()}, which keeps IDs opaque and stable across the
API boundary; the sole exception is \texttt{auth.users.id}, a native
\texttt{uuid} that \texttt{candidates.user\_id} references. The
\texttt{updated\_at} column is maintained by a PostgreSQL trigger rather than
in application code.

\paragraph{Access control without RLS.}
Row-Level Security is deliberately \emph{disabled} on the application tables.
Every database call is made server-side by code holding the service-role
key, and authorisation is enforced one layer out --- in the middleware and
in route helpers such as \texttt{require-hr.ts} --- rather than in the
database. This concentrates the authorisation logic in one auditable place,
at the explicit cost of making the server-side trust boundary critical;
that cost is precisely why the authentication and route-gating tests of
\chapref{cap:validation} are central to the design.

\paragraph{Relationships and integrity.}
The schema is centred on the \texttt{candidates} table, from which radiate
one-to-many relationships to experiences, education, languages, skills, tags,
notes, assessments, improvement tracks, job matches, applications and
notifications, plus a one-to-one notification-preference row. The
\texttt{jobs} table similarly fans out to its applications, cached matches,
assessment templates and shortlist entries. Junction-style tables with
composite uniqueness express the genuinely many-to-many relationships ---
\texttt{job\_shortlists} is unique on \texttt{(job\_id, candidate\_id)},
\texttt{candidate\_languages} on \texttt{(candidate\_id, language)}, and
\texttt{job\_applications} on \texttt{(job\_id, candidate\_id)} --- so the
database, not the application, guarantees that a candidate cannot apply to
the same job twice or hold two records for the same language. Foreign keys
are declared throughout, both for referential integrity and because the
Supabase PostgREST layer requires a real foreign key to resolve its
relational join syntax.

\paragraph{Indexing strategy.}
Indexes target the queries the application actually issues. The
\texttt{candidates} table is indexed on \texttt{email}, \texttt{status},
\texttt{overall\_cv\_score} and \texttt{country}, which are the columns the
talent-pool filters sort and restrict on. The array column
\texttt{experiences.fields\_of\_work} carries a \textsc{gin} index so that the
matching engine can test field membership efficiently across the whole
experience corpus. The \texttt{hr\_dashboard\_widgets} table is indexed on
\texttt{(user\_id, position)} to serve each recruiter's ordered set of custom
charts in a single scan.

\figref{fig:er} shows the core of the resulting schema. The two darker
entities, \texttt{candidates} and \texttt{jobs}, are the aggregates from
which everything else fans out; the lighter join entities
(\texttt{job\_applications}, \texttt{job\_matches}, \texttt{job\_shortlists})
are exactly the genuinely many-to-many relationships protected by composite
uniqueness. Attributes are omitted here for legibility and are listed in full
in \appref{app:db-schema}.

\begin{figure}[ht]
\centering
\resizebox{\textwidth}{!}{%
\begin{tikzpicture}[font=\scriptsize,
  ent/.style={draw, fill=lstkeyword!15, rounded corners, align=center, minimum height=0.6cm, inner sep=3pt},
  agg/.style={draw, fill=lstkeyword!38, rounded corners, align=center, minimum height=0.7cm, inner sep=3pt, font=\scriptsize\bfseries},
  jent/.style={draw, fill=lstcomment!18, rounded corners, align=center, minimum height=0.6cm, inner sep=3pt},
  rl/.style={draw, lstcomment}]
\node[agg] (cand) at (0,0) {candidates};
\node[ent] (exp) at (-4.6,2.4) {experiences};
\node[ent] (edu) at (-4.6,1.0) {education};
\node[ent] (lan) at (-4.6,-0.4) {candidate\_languages};
\node[ent] (ski) at (-4.6,-1.8) {skills};
\node[ent] (isess) at (-3.4,-3.2) {interview\_sessions};
\node[jent] (japp) at (3.6,2.6) {job\_applications};
\node[jent] (jmat) at (3.6,1.2) {job\_matches};
\node[jent] (jsho) at (3.6,-0.2) {job\_shortlists};
\node[agg] (job) at (7.4,1.2) {jobs};
\node[ent] (ass) at (3.6,-2.0) {assessments};
\node[ent] (ares) at (7.4,-2.0) {assessment\_results};
\foreach \c in {exp,edu,lan,ski,isess,japp,jmat,jsho,ass} \draw[rl] (cand) -- node[pos=0.82,font=\tiny]{$*$} (\c);
\foreach \c in {japp,jmat,jsho,ass} \draw[rl] (job) -- node[pos=0.82,font=\tiny]{$*$} (\c);
\draw[rl] (ass) -- node[pos=0.82,font=\tiny]{$1$} (ares);
\end{tikzpicture}}
\caption[Core entity--relationship diagram]{Core of the relational schema
(attributes omitted); ``$*$'' marks the many side of each relationship.}
\label{fig:er}
\end{figure}

%----------------------------------------------------------------------------------------

\section{Job-Anchored Matching Algorithm}
\label{sec:design-matching}

A central design decision is that there is \emph{no} universal
``candidate score'' for hiring; matching is always candidate~$\times$~a
specific job. The system therefore exposes two distinct numbers with
distinct meanings:

\begin{description}
	\item[Quality.] A CV-intrinsic profile score
	(\texttt{overall\_cv\_score}) reflecting completeness, education,
	languages and location. It is independent of any job and is useful only
	as a prefilter, never as a hiring signal.
	\item[Fit.] Computed live for a chosen job by \texttt{computeJobFit}.
\end{description}

\texttt{computeJobFit} is a pure function evaluating seven criteria: field
of work, experience in that field, seniority, required skills, preferred
skills, languages and education. The overall fit is the average of the
\emph{applicable} criteria only --- a criterion the job does not specify is
excluded rather than counted as zero, so jobs are not penalised for being
unspecific. A separate \texttt{isEligible} flag is the logical conjunction
of the applicable ``met'' flags, distinguishing ``ranked but not strictly
eligible'' from ``eligible''. Each criterion also returns its own explanation,
so the final ranking is fully traceable rather than a single opaque number.

Two properties of this design are worth stating precisely, because they bound
what the ranking may and may not claim. Because \texttt{computeJobFit} is a
pure function, its output is deterministic and reproducible; that consistency,
together with the per-criterion correctness of each rule, is demonstrated by
the unit suite in \secref{sec:val-unit}. What the design does \emph{not}
establish is that the resulting order coincides with an expert recruiter's own
ranking of the same candidates: that agreement was left unvalidated for want of
real application data and recruiter time, and is recorded honestly as a
limitation in \secref{sec:val-limitations}.

\figref{fig:match-flow} traces the orchestration end to end. The element
worth drawing out is the output cache at the final step: the ranked result is
persisted to \texttt{job\_matches}, so a recruiter reopening the screen reads the
stored ranking rather than triggering a full re-score.

\begin{figure}[ht]
\centering
\begin{tikzpicture}[font=\scriptsize, node distance=6mm,
  flowstep/.style={draw, rounded corners, fill=lstkeyword!18, text width=8.0cm, align=center, minimum height=0.65cm},
  arr/.style={-{Stealth}, draw}]
\node[flowstep] (s1) {Lazy-parse job description $\rightarrow$ versioned requirements (cached until source or schema version changes)};
\node[flowstep, below=of s1] (s2) {Load candidates with experiences, languages, education and skills};
\node[flowstep, below=of s2] (s3) {Build per-field experience vector from \texttt{fields\_of\_work}};
\node[flowstep, below=of s3] (s4) {For each candidate: \texttt{computeJobFit} over 7 criteria (pure function)};
\node[flowstep, below=of s4] (s5) {Average applicable criteria $\rightarrow$ fit \%; AND of ``met'' $\rightarrow$ \texttt{isEligible}};
\node[flowstep, below=of s5] (s6) {Persist top-100 to \texttt{job\_matches} cache; render ranked screen};
\foreach \a/\b in {s1/s2,s2/s3,s3/s4,s4/s5,s5/s6} \draw[arr] (\a) -- (\b);
\end{tikzpicture}
\caption[Job-matching orchestration]{Job-matching orchestration, end to end.}
\label{fig:match-flow}
\end{figure}

%----------------------------------------------------------------------------------------

\section{AI Pipelines}
\label{sec:design-ai}

\subsection{CV Parsing Pipeline}
\label{sec:design-cv}

At the design level the CV pipeline is organised around a single principle:
the LLM is an \emph{untrusted boundary} between an unstructured document and a
typed Candidate aggregate. The file is stored, its text extracted, and that
text sent to the primary model under a structured-output
prompt; the response is validated against a strict Zod schema before any of it
reaches the database, with an automatic retry on a parse failure and a
fallback to the secondary provider if the primary itself fails. The concrete stage-by-stage
breakdown, the deduplication tiers and the per-extraction telemetry are
detailed in \chapref{cap:implementation}.

\figref{fig:seq-cv} renders this pipeline as an interaction sequence. The
two self-messages on the use case --- text extraction and schema validation
with its single retry --- are where the untrusted-boundary discipline of
\secref{sec:sota-structured} takes concrete form in the design. The full
structured-output prompt that drives the extraction is reproduced in
\appref{app:prompts}.

\begin{figure}[H]
\centering
\resizebox{\textwidth}{!}{%
\begin{tikzpicture}[font=\scriptsize, >=Stealth,
  part/.style={draw, fill=lstkeyword!15, align=center, text width=2.0cm, minimum height=0.7cm},
  m/.style={-{Stealth}, draw},
  r/.style={-{Stealth}, draw, dashed, lstcomment}]
\node[part] (br) at (0,0) {Browser};
\node[part] (uc) at (3.4,0) {Upload use case};
\node[part] (st) at (6.2,0) {Supabase Storage};
\node[part] (llm) at (9.0,0) {LLM Groq/OpenAI};
\node[part] (db) at (11.8,0) {Database};
\foreach \px in {0,3.4,6.2,9.0,11.8} \draw[lstcomment!50,dashed] (\px,-0.45) -- (\px,-8.4);
\draw[m] (0,-0.9) -- node[above]{upload CV (PDF/DOCX)} (3.4,-0.9);
\draw[m] (3.4,-1.7) -- node[above]{store file} (6.2,-1.7);
\draw[r] (6.2,-2.4) -- node[above]{file URL} (3.4,-2.4);
\draw[m] (3.4,-3.1) -- ++(1.2,0) -- ++(0,-0.5) -- (3.4,-3.6);
\node[anchor=west,font=\tiny] at (4.65,-3.35) {extract text};
\draw[m] (3.4,-4.3) -- node[above]{structured-output prompt} (9.0,-4.3);
\draw[r] (9.0,-5.0) -- node[above]{JSON candidate (retry on fail $\rightarrow$ OpenAI)} (3.4,-5.0);
\draw[m] (3.4,-5.7) -- ++(1.2,0) -- ++(0,-0.5) -- (3.4,-6.2);
\node[anchor=west,font=\tiny] at (4.65,-5.95) {validate (Zod) + retry};
\draw[m] (3.4,-6.9) -- node[above]{persist candidate + children} (11.8,-6.9);
\draw[r] (11.8,-7.5) -- node[above]{ok} (3.4,-7.5);
\draw[r] (3.4,-8.1) -- node[above]{parsed candidate} (0,-8.1);
\end{tikzpicture}}
\caption[CV-parsing sequence]{The CV-parsing pipeline as an interaction
sequence.}
\label{fig:seq-cv}
\end{figure}

\subsection{Job Description Extractor}
\label{sec:design-jd}

Job descriptions are parsed by the same structured-output approach through a
dedicated service, into a \emph{versioned} requirements schema. Parsing is
\textbf{lazy}: a job description is parsed the first time an HR user opens its
ranking screen, not in bulk, and the result is cached with its schema version
so it is reused until either the source text or the schema version changes.
This keeps LLM usage proportional to actual use. The extractor's prompt ---
which normalises multilingual education and language terms and infers
requirements from duties-only postings --- is reproduced in
\appref{app:prompts}.

\subsection{AI Interviewer (FastAPI Sidecar)}
\label{sec:design-interviewer}

The conversational assessment runs in a separate Python FastAPI service (the
second member's module; \secref{sec:team}). The design point relevant here is
the \emph{contract}. The browser opens an interview popup that talks to the
Next.js API routes; those routes create a session, issue a short-lived
HMAC-SHA256 token (ten-minute TTL, with only its hash stored), proxy the
conversation turns to the sidecar, and persist the final evaluation. This
per-session token governs the live interview only and is distinct from the
48-hour \texttt{MAGIC\_LINK\_EXPIRY\_HOURS} window (\tabref{tab:value-objects})
that keeps an assessment-invitation link valid until the candidate starts. The
interview runs in one of two modes (technical or language), and the sidecar's
evaluator returns a structured result that the Next.js side persists as JSONB;
the conversational internals and the mode prompts are the second member's,
summarised in \appref{app:prompts}.

\figref{fig:seq-interview} traces the token and session flow across the
boundary. The Next.js core owns identity and persistence; the sidecar owns the
conversation. The only thing crossing between them is a short-lived,
hash-verified token and the structured turns it authorises --- a contract
small enough to reason about in full.

\begin{figure}[H]
\centering
\resizebox{\textwidth}{!}{%
\begin{tikzpicture}[font=\scriptsize, >=Stealth,
  part/.style={draw, fill=lstkeyword!15, align=center, text width=2.0cm, minimum height=0.7cm},
  m/.style={-{Stealth}, draw},
  r/.style={-{Stealth}, draw, dashed, lstcomment}]
\node[part] (cand) at (0,0) {Candidate browser};
\node[part] (api) at (3.8,0) {Next.js API};
\node[part] (side) at (7.6,0) {FastAPI sidecar};
\node[part] (db) at (10.8,0) {Database};
\foreach \px in {0,3.8,7.6,10.8} \draw[lstcomment!50,dashed] (\px,-0.45) -- (\px,-10.0);
\draw[m] (0,-0.9) -- node[above]{start interview} (3.8,-0.9);
\draw[m] (3.8,-1.7) -- node[above]{create session + store token hash} (10.8,-1.7);
\draw[r] (3.8,-2.4) -- node[above]{signed HMAC token (10-min TTL)} (0,-2.4);
\draw[m] (0,-3.1) -- node[above]{turn (answer) + token} (3.8,-3.1);
\draw[m] (3.8,-3.8) -- ++(1.2,0) -- ++(0,-0.5) -- (3.8,-4.3);
\node[anchor=west,font=\tiny] at (5.05,-4.05) {verify token hash};
\draw[m] (3.8,-5.0) -- node[above]{proxy turn} (7.6,-5.0);
\draw[r] (7.6,-5.7) -- node[above]{next question} (3.8,-5.7);
\draw[r] (3.8,-6.4) -- node[above]{next question} (0,-6.4);
\draw[m] (0,-7.1) -- node[above]{finish + token} (3.8,-7.1);
\draw[m] (3.8,-7.8) -- node[above]{evaluate} (7.6,-7.8);
\draw[r] (7.6,-8.5) -- node[above]{CEFR / technical result} (3.8,-8.5);
\draw[m] (3.8,-9.1) -- node[above]{persist result (JSONB)} (10.8,-9.1);
\draw[r] (3.8,-9.7) -- node[above]{result} (0,-9.7);
\end{tikzpicture}}
\caption[Interview token and session sequence]{Token and session flow between
the Next.js core and the FastAPI interviewer sidecar.}
\label{fig:seq-interview}
\end{figure}

%----------------------------------------------------------------------------------------

\section{Security and Compliance}
\label{sec:design-security}

\subsection{Authentication and Authorisation}
\label{sec:design-authz}

Authentication is delegated entirely to Supabase Auth using Google OAuth.
The user's role --- \texttt{candidate} or \texttt{hr} --- is stored in
\texttt{app\_metadata.role}, set server-side on first sign-in with the
service-role key so that it cannot be altered from the client. The Next.js
middleware refreshes the session on every request and gates access: an
unauthenticated caller to an API route receives \texttt{401}, a non-HR
caller to an HR-only prefix receives \texttt{403}, and dashboard routes are
protected accordingly. Authorisation thus lives in the application edge,
consistent with the decision to disable database RLS.

\subsection{GDPR Alignment}
\label{sec:design-gdpr}

The data model is deletion-aware to support the six-month retention rule:
candidate data can be removed through standard operations, and only the
inferences actually needed are stored. A per-candidate interaction history
provides an audit trail of status changes, emails and campaign outreach, and
the candidate self-service portal gives data subjects visibility over their
own record. The rights of erasure and access are exercised directly by the
data subject rather than through a back-office request --- the self-service
controls that erase the full record and its stored documents, or export the
profile, are described in \secref{sec:impl-engagement}. The remaining gap --- automating
the time-based purge so that retention is enforced without a manual trigger
--- is recorded as future work in \chapref{cap:conclusions}.

\subsection{Hardening and Transparency}
\label{sec:design-hardening}

Beyond authorisation, a handful of measures address the common web
risk classes of the OWASP Top~10 \parencite{OWASPTop10}. Injection is closed at the data edge by
escaping user search terms before they reach the PostgREST filter
(\chapref{cap:implementation}); broken access control is mitigated by the
edge gate above and by setting roles server-side; and the interviewer
boundary is protected by short-lived, hash-verified HMAC tokens
(\figref{fig:seq-interview}) so that a leaked token neither persists nor
reveals a usable secret. Transport and storage encryption (TLS in transit,
AES-256 at rest) are inherited from the managed platform, the service-role key
never leaves the server, and the structured logger redacts secret-bearing
keys so credentials cannot surface in logs.

Transparency is treated as a first-class concern rather than an afterthought.
A public privacy policy --- reachable before authentication --- names the data
controller, the lawful basis for each processing purpose, the sub-processors
and international-transfer safeguards, and the full set of data-subject rights.
Critically for an AI-assisted tool, it discloses the automated-profiling
position under \gls{gdpr} Article~22: the match and assessment scores
\emph{assist} a human reviewer rather than decide automatically, and a manual
review can be requested. A first-visit consent banner records the visitor's
choice and links to that policy; because the platform sets only
strictly-necessary session cookies and no tracking or advertising cookies,
the banner is informational rather than a processing gate --- a position that
would tighten into an explicit opt-in the moment analytics cookies were
introduced.

%----------------------------------------------------------------------------------------

\section{Frontend Design}
\label{sec:design-frontend}

The interface is built with shadcn/ui components over Tailwind CSS, which
gives a consistent, accessible primitive set without committing to a heavy
component framework. The information architecture is split by persona: the
\textbf{HR dashboard} is organised around the recruiter's workflow (talent
pool, jobs and ranking, shortlists, assessments, campaigns, notifications,
interaction history and analytics), while the \textbf{candidate portal}
exposes only self-service actions (registration, profile and notification
preferences, document upload and editing, applications to jobs, internships
and the ambassador program, assessments, application-status tracking and
account deletion). The same application serves both, with navigation and
permitted actions driven by the authenticated role, so a single codebase
presents two coherent experiences.

%----------------------------------------------------------------------------------------

\section{Design Decision Log}
\label{sec:design-decisions}

\tabref{tab:design-decisions} records the principal design decisions, the
alternatives weighed, and the trade-off accepted in each case.

\begin{table}[t]
\centering
\caption[Design decision log]{Principal design decisions and their trade-offs.}
\label{tab:design-decisions}
\small
\begin{tabular}{p{0.24\linewidth}p{0.26\linewidth}p{0.38\linewidth}}
\toprule
\textbf{Decision} & \textbf{Alternative} & \textbf{Rationale / trade-off} \\
\midrule
Supabase (DB+Auth+Storage) & ORM + separate Neon, Blob, auth & One managed provider; fewer integration layers; migration cost paid once \\
No ORM, raw client & Prisma/Drizzle & Full control + transparency; cost is manual key conversion \\
RLS disabled, app-side authz & Per-row RLS policies & Single auditable authz layer; server trust boundary becomes critical \\
Rule-based \texttt{computeJobFit} & Embedding similarity & Per-criterion explainability over opaque ranking \\
Groq primary, OpenAI fallback & Single provider & Outage becomes degradation, not failure \\
FastAPI sidecar for interviewer & In-process JS & Clean language/process boundary; parallel development \\
Text UUID primary keys & DB serial / native uuid & Opaque, collision-resistant, distributed-safe IDs \\
\bottomrule
\end{tabular}
\end{table}
